% Artykuł własny: LNCS, po polsku, IMRaD, max 6 stron bez bibliografii.
% Kompilacja: latexmk -pdf main.tex
% Tabele w tables/ generuje research/summarize.py -- nie edytować ręcznie.
% \TODO{...} oznacza miejsca do uzupełnienia przez grupę.
\documentclass[runningheads]{llncs}
\usepackage[T1]{fontenc}
\usepackage[utf8]{inputenc}
\usepackage{polski}
\usepackage{lmodern}
\usepackage{amsmath}
\usepackage{booktabs}
\usepackage{graphicx}
\usepackage{xcolor}
\usepackage{hyperref}
\newcommand{\TODO}[1]{\textcolor{red}{[TODO: #1]}}
\renewcommand{\abstractname}{Streszczenie}
\renewcommand{\keywordname}{\textbf{Słowa kluczowe:}}
\renewcommand{\tablename}{Tabela}
\renewcommand{\figurename}{Rys.}
\renewcommand{\refname}{Bibliografia}

\begin{document}

\title{Ocena silników OCR i modeli wizyjno-językowych w~rozpoznawaniu tekstu polskiego na danych syntetycznych i~rzeczywistych}
\titlerunning{OCR i VLM dla tekstu polskiego}
\author{\TODO{Autor 1} \and \TODO{Autor 2} \and \TODO{Autor 3}}
\authorrunning{\TODO{Autorzy}}
\institute{\TODO{Uczelnia, wydział}\\\email{\TODO{e-mail}}}

\maketitle

\begin{abstract}
Systemy OCR i~modele wizyjno-językowe (VLM) trenowane są głównie na danych angielskich, przez co
gorzej radzą sobie ze znakami diakrytycznymi innych języków. Przenosimy protokół badania
Yılmaza i~in.~\cite{yilmaz2026} z~języka tureckiego na polski: generujemy syntetyczny zbiór 2200 obrazów
z~\emph{Lalki} B.~Prusa (tekst drukowany i~odręczny), kontrolując obecność polskich znaków, długość słowa,
kontekst zdania i~pięć rodzajów zniekształceń, oraz oceniamy modele na 1968 skanach odręcznie przepisanej
Inwokacji z~\emph{Pana Tadeusza}. Porównujemy siedem otwartych systemów: Qwen2.5-VL 7B, InternVL3-8B,
DeepSeek-VL2, GOT-OCR2, PaddleOCR oraz Tesseract i~EasyOCR z~modelami polskimi. Najlepszy jest Qwen2.5-VL
(CER 0,059 na danych syntetycznych i~0,102 na skanach). Ranking na syntetycznym piśmie odręcznym zgadza
się z~rankingiem na skanach (Spearman $\rho=0{,}89$, $p=0{,}007$). Pokazujemy też, że słownik znaków
PaddleOCR dla języków łacińskich nie zawiera liter ń, ś, ź, ż. Kod i~dane udostępniamy publicznie.
\keywords{OCR \and modele wizyjno-językowe \and język polski \and dane syntetyczne \and znaki diakrytyczne}
\end{abstract}

%------------------------------------------------------------------------------
\section{Wprowadzenie}

Rozpoznawanie tekstu przeszło w~ostatnich latach drogę od klasycznych silników OCR do dużych
modeli wizyjno-językowych (VLM), które odczytują tekst z~obrazu bez osobnych etapów detekcji
i~segmentacji. Modele te uczone są głównie na danych angielskich i~chińskich.

Yılmaz i~in.~\cite{yilmaz2026} zaproponowali syntetyczny benchmark dla języka tureckiego i~pokazali,
że znaki specyficzne dla języka są główną trudnością dla większości modeli. Język polski ma
dziewięć liter z~diakrytykami (ą, ć, ę, ł, ń, ó, ś, ź, ż), z~których część różni się od liter
łacińskich jedynie drobnym znakiem (ogonek, kreska, kropka).

\textbf{Cel.} Celem artykułu jest ocena, jak dokładnie otwarte silniki OCR i~modele VLM rozpoznają
tekst polski oraz jak na wynik wpływają polskie znaki, długość słowa, kontekst zdania, rodzaj pisma
i~zniekształcenia obrazu. Pytania badawcze: (1) który model najlepiej rozpoznaje polski tekst,
(2) jak na błąd wpływają wymienione czynniki oraz (3) czy ranking modeli na danych syntetycznych
zgadza się z~rankingiem na prawdziwym piśmie odręcznym.

\textbf{Elementy nowości:}
(1) przeniesienie protokołu~\cite{yilmaz2026} na język polski wraz z~otwartym generatorem;
(2) uwzględnienie klasycznych silników z~modelem polskim (Tesseract \texttt{pol}, EasyOCR \texttt{pl}),
których autorzy nie badali;
(3) zbiór rzeczywisty: 1968 skanów odręcznie przepisanej Inwokacji, na którym sprawdzamy, czy wyniki
syntetyczne przekładają się na prawdziwe pismo;
(4) analiza błędów wyłącznie diakrytycznych i~ograniczeń słowników znaków;
(5) weryfikacja odtwarzalności badania bazowego.

%------------------------------------------------------------------------------
\section{Przegląd literatury}

Tesseract~\cite{smith2007} i~PP-OCR~\cite{li2022ppocrv3} to klasyczne systemy z~osobnymi modułami detekcji
i~rozpoznawania, a~EasyOCR~\cite{easyocr} łączy detektor z~rozpoznawaniem sekwencyjnym. Nowsze modele
rozpoznają tekst jednym przebiegiem: GOT-OCR2~\cite{wei2024got} jest modelem wyspecjalizowanym w~OCR,
a~Qwen2.5-VL~\cite{bai2025qwen25vl}, InternVL3~\cite{zhu2025internvl3} i~DeepSeek-VL2~\cite{wu2024deepseekvl2}
to ogólne modele wizyjno-językowe. Dane syntetyczne generowane biblioteką TRDG~\cite{trdg} pozwalają
kontrolować właściwości tekstu; wykorzystują je Yılmaz i~in.~\cite{yilmaz2026}, których protokół
powtarzamy. \TODO{Prace o~OCR dla języka polskiego i~zbiór PHSF (max 10 pozycji łącznie w~bibliografii).}

%------------------------------------------------------------------------------
\section{Materiał i metody}

\subsection{Zbiór syntetyczny}
Tekst źródłowy to \emph{Lalka}, tom pierwszy i~drugi~\cite{wolnelektury} (domena publiczna). Pulę słów
stanowią wyrazy co najmniej trzyliterowe występujące w~tekście co najmniej dwa razy, a~pulę zdań --
zdania liczące 4--8 słów, złożone wyłącznie z~liter i~interpunkcji.
Obrazy generujemy biblioteką TRDG 1.8.0~\cite{trdg} z~parametrami autorów badania bazowego: wysokość
64~px, marginesy 15~px, czarny tekst na białym tle. Tekst drukowany renderujemy 15 fontami, a~odręczny
19 popularnymi fontami odręcznymi z~Google Fonts. Każdy font sprawdzamy w~tabeli \texttt{cmap}
pod kątem wszystkich znaków korpusu, aby wykluczyć brakujące glify.

Dla każdego rodzaju pisma (drukowane, odręczne) tworzymy cztery kategorie:
(1) słowa z~polskimi znakami i~bez nich,
(2) słowa krótkie (${<}10$ znaków) i~długie (${>}10$),
(3) pojedyncze słowo i~zdanie,
(4) zniekształcenia: obraz czysty, rozmycie (promień 2), obrót (do $\pm16^\circ$), szum gaussowski
($\sigma=50$) i~niska rozdzielczość (wysokość 40~px), stosowane do tych samych słów.
Daje to 22 komórki po 100 obrazów, łącznie 2200 obrazów. Losowanie jest deterministyczne (ziarno 42).

\subsection{Zbiór rzeczywisty (Gold Set)}
Używamy zbioru PHSF \TODO{cytowanie PHSF}: skanów stron, na których respondenci odręcznie przepisali
cztery pierwsze wersy Inwokacji z~\emph{Pana Tadeusza}. Bierzemy wszystkie 1968 skanów (szarość,
szerokość 2500--5900~px, 4 linie). Zmniejszamy je do 2048~px dłuższego boku, tak samo dla wszystkich
modeli, i~przetwarzamy w~całości, bez cięcia na linie. Tekst wzorcowy to Inwokacja po tej samej
normalizacji co predykcje (małe litery, bez interpunkcji); błędy przepisujących pozostają w~danych.

\subsection{Modele}
Wszystkie modele są otwarte i~uruchamiane w~Google Colab (GPU NVIDIA A100), każdy we własnym
środowisku z~przypiętymi wersjami pakietów i~commitem modelu. Prompt VLM to prompt autorów
z~,,Turkish'' zamienionym na ,,Polish''.
\begin{itemize}
  \item \textbf{Qwen2.5-VL 7B}~\cite{bai2025qwen25vl} (limit 512 tokenów wyjścia),
    \textbf{InternVL3-8B}~\cite{zhu2025internvl3}, \textbf{DeepSeek-VL2-tiny}~\cite{wu2024deepseekvl2};
  \item \textbf{GOT-OCR2}~\cite{wei2024got};
  \item \textbf{PaddleOCR} 2.9.1~\cite{li2022ppocrv3}, model \texttt{pl};
  \item \textbf{Tesseract}~\cite{smith2007}, model \texttt{pol}: tryb jednej linii (\texttt{--psm 7})
    dla obrazów syntetycznych, tryb bloku tekstu (\texttt{--psm 6}) dla skanów;
  \item \textbf{EasyOCR}~\cite{easyocr}, model \texttt{pl}.
\end{itemize}
Pominęliśmy TrOCR (checkpoint \texttt{large-printed} rozpoznaje pojedyncze linie angielskiego druku, więc
nie obsługuje całych stron) oraz modele zamknięte (GPT-4o, Claude). Phi-3.5-vision i~Pixtral-12B uruchomiliśmy,
ale przebiegu nie dokończyliśmy, bo Colab przestał przydzielać odpowiednie GPU (Phi: 833 z~2200 obrazów
syntetycznych; Pixtral: 200 z~1968 skanów), więc nie wchodzą do analizy.

\subsection{Miary}
Niech $t_i$ oznacza tekst wzorcowy, a~$p_i$ predykcję modelu dla obrazu $i$. Oba teksty normalizujemy:
Unicode NFC, małe litery, \emph{usunięta interpunkcja}, scalone białe znaki; polskie znaki zostają.
Interpunkcję usuwamy, bo przepisujący stosują ją niespójnie, a~modele zwracają ją losowo, więc
mierzylibyśmy nawyki, a~nie rozpoznawanie liter. Główną miarą jest współczynnik błędów znakowych
\begin{equation}\label{eq:cer}
  \mathrm{CER}_i = \min\left(1,\ \frac{S_i + D_i + I_i}{N_i}\right),
\end{equation}
gdzie $S_i$, $D_i$, $I_i$ to liczby podstawień, usunięć i~wstawień znaków w~minimalnym ciągu edycji,
a~$N_i = |t_i|$ to liczba znaków wzorca. Obcięcie do~1 odpowiada protokołowi~\cite{yilmaz2026}.
Współczynnik błędów słów $\mathrm{WER}_i$ liczymy tym samym wzorem na ciągach słów i~podajemy tylko dla
zdań i~skanów, bo dla pojedynczego słowa przyjmuje on wartości 0 lub~1.

Wynik dla zbioru obrazów $A$ (komórki, rodzaju pisma lub całego zbioru) to średnia
\begin{equation}\label{eq:mean}
  \overline{\mathrm{CER}}(A) = \frac{1}{|A|}\sum_{i\in A}\mathrm{CER}_i .
\end{equation}
Wpływ polskich znaków opisujemy odsetkiem \emph{błędów wyłącznie diakrytycznych}: wśród błędnych
predykcji $E=\{i: p_i\neq t_i\}$ liczymy te, które stają się poprawne po usunięciu diakrytyków
funkcją $s(\cdot)$ (ą$\to$a, ł$\to$l itd.):
\begin{equation}\label{eq:diacr}
  R_{\mathrm{diakr}} = \frac{|\{i\in E: s(p_i) = s(t_i)\}|}{|E|}.
\end{equation}
Na skanach podajemy ponadto odsetek polskich liter wzorca, które model odczytał błędnie (zamienione
lub pominięte).

\subsection{Analiza statystyczna}
Niepewność średniej~\eqref{eq:mean} szacujemy jak w~\cite{yilmaz2026} nieparametrycznym bootstrapem:
z~wartości $\mathrm{CER}_i$ losujemy ze zwracaniem $B=1000$ prób o~liczności $|A|$, a~95\% przedział
ufności wyznaczają percentyle 2,5 i~97,5 rozkładu średnich. Wpływ czynnika mierzymy różnicą średnich
między grupą z~czynnikiem $A$ i~grupą odniesienia $B$:
\begin{equation}\label{eq:delta}
  \Delta = \overline{\mathrm{CER}}(A) - \overline{\mathrm{CER}}(B),
\end{equation}
np.\ $A$ -- słowa z~polskimi znakami, $B$ -- słowa bez nich; dla zniekształceń $B$ to obrazy czyste.
Przedział ufności dla $\Delta$ wyznaczamy, losując niezależnie obie grupy; efekt uznajemy za istotny,
gdy przedział nie zawiera zera. Zgodność rankingów na danych syntetycznych i~rzeczywistych oceniamy
współczynnikiem korelacji rang Spearmana między średnim CER na zdaniach pisma odręcznego (obrazy
czyste) a~średnim CER na skanach.

%------------------------------------------------------------------------------
\section{Wyniki}

\subsection{Zbiór syntetyczny}
Tabela~\ref{tab:overall} przedstawia wyniki zbiorcze (wyniki w~komórkach są w~repozytorium).
Najniższy średni CER ma Qwen2.5-VL 7B (0,059), następnie DeepSeek-VL2 (0,143) i~InternVL3-8B (0,147);
przedziały ufności tych dwóch ostatnich się pokrywają. Za nimi są EasyOCR (0,174), PaddleOCR (0,183),
Tesseract (0,206) i~GOT-OCR2 (0,320). Pismo odręczne jest trudniejsze dla wszystkich modeli
($\Delta$ od $+0{,}08$ u~Qwena do $+0{,}24$ u~Tesseracta). Na czystym druku Tesseract i~EasyOCR są niemal
bezbłędne (CER ${\le}0{,}002$), ale tracą dużo przy szumie i~niskiej rozdzielczości.

\begin{table}
  \caption{Średni CER~\eqref{eq:mean} według rodzaju pisma z~95\% przedziałem ufności (bootstrap), WER dla
  zdań, odsetek błędów wyłącznie diakrytycznych~\eqref{eq:diacr} i~czas na obraz (Colab A100).}
  \label{tab:overall}
  \centering\small
  \begin{tabular}{lrrrcrrr}
\toprule
Model & \multicolumn{4}{c}{CER} & WER & Diakr. & s/obraz \\
 & druk & pismo & razem & 95\% CI & zdania & [\%] & \\
\midrule
Qwen2.5-VL 7B & .019 & .099 & .059 & [.053, .066] & .041 & 31 & 0.25 \\
InternVL3-8B & .093 & .200 & .147 & [.138, .155] & .244 & 26 & 0.57 \\
DeepSeek-VL2 & .100 & .185 & .143 & [.135, .151] & .351 & 33 & 0.28 \\
Tesseract & .085 & .326 & .206 & [.192, .220] & .215 & 3 & 0.12 \\
EasyOCR & .077 & .271 & .174 & [.162, .186] & .290 & 10 & 0.02 \\
PaddleOCR & .116 & .251 & .183 & [.173, .194] & .452 & 29 & 0.03 \\
GOT-OCR2 & .245 & .395 & .320 & [.310, .331] & .747 & 10 & 0.27 \\
\bottomrule
\end{tabular}

\end{table}

\begin{table}
  \caption{Wpływ czynników: różnica średniego CER $\Delta$~\eqref{eq:delta} między grupą z~czynnikiem
  a~grupą odniesienia (dla zniekształceń: obrazy czyste). Pogrubienie: 95\% przedział ufności nie
  zawiera zera.}
  \label{tab:effects}
  \centering\small
  \begin{tabular}{lrrrrrrr}
\toprule
Czynnik & Qwen2.5-VL 7B & InternVL3-8B & DeepSeek-VL2 & Tesseract & EasyOCR & PaddleOCR & GOT-OCR2 \\
\midrule
polskie znaki & \textbf{+.046} & \textbf{+.139} & \textbf{+.180} & +.015 & \textbf{+.029} & \textbf{+.159} & \textbf{+.227} \\
długie słowa & -.009 & \textbf{-.044} & \textbf{-.040} & -.004 & -.011 & \textbf{-.043} & -.000 \\
zdanie zamiast słowa & \textbf{-.032} & \textbf{-.043} & -.016 & +.008 & \textbf{+.049} & \textbf{-.057} & \textbf{-.067} \\
pismo odręczne & \textbf{+.080} & \textbf{+.108} & \textbf{+.085} & \textbf{+.241} & \textbf{+.194} & \textbf{+.136} & \textbf{+.150} \\
rozmycie & \textbf{+.084} & \textbf{+.183} & \textbf{+.161} & \textbf{+.239} & \textbf{+.435} & \textbf{+.052} & \textbf{+.101} \\
obrót & +.017 & +.004 & -.003 & \textbf{+.118} & \textbf{+.038} & -.000 & +.013 \\
szum & +.003 & \textbf{+.059} & +.004 & \textbf{+.522} & +.010 & +.014 & +.042 \\
niska rozdzielczość & \textbf{+.205} & \textbf{+.234} & \textbf{+.256} & \textbf{+.535} & \textbf{+.683} & \textbf{+.407} & \textbf{+.186} \\
\bottomrule
\end{tabular}

\end{table}

Tabela~\ref{tab:effects} pokazuje wpływ poszczególnych czynników.
\textbf{Polskie znaki} zwiększają błąd najbardziej u~GOT-OCR2 ($+0{,}23$), DeepSeek-VL2 ($+0{,}18$),
PaddleOCR ($+0{,}16$) i~InternVL3 ($+0{,}14$); u~Qwena jest to $+0{,}05$, a~u~Tesseracta i~EasyOCR
niewiele ($+0{,}02$ i~$+0{,}03$, u~Tesseracta nieistotnie).
\textbf{Długość słowa} nie pogarsza wyników: efekt jest zerowy lub lekko ujemny, a~u~InternVL3, DeepSeek-VL2
i~PaddleOCR istotnie ujemny (ok.\ $-0{,}04$), bo dłuższe słowa mają więcej kontekstu.
\textbf{Zdanie zamiast słowa} obniża błąd u~większości modeli (PaddleOCR $-0{,}06$, GOT-OCR2 $-0{,}07$,
InternVL3 $-0{,}04$); wyjątkiem jest EasyOCR ($+0{,}05$).
\textbf{Zniekształcenia:} niska rozdzielczość psuje wszystkie modele, ale najbardziej klasyczne
(EasyOCR $+0{,}68$, Tesseract $+0{,}53$, PaddleOCR $+0{,}41$) wobec $+0{,}19$--$+0{,}26$ u~pozostałych;
szum zrywa niemal wyłącznie Tesseracta ($+0{,}52$); rozmycie najbardziej dotyka EasyOCR ($+0{,}44$)
i~Tesseracta ($+0{,}24$); obrót ma niewielki wpływ poza Tesseractem ($+0{,}12$).

\subsection{Skany odręcznej Inwokacji}
Na skanach (tab.~\ref{tab:gold}) błędy są znacznie większe niż na obrazach syntetycznych, bo
modele czytają całą stronę z~pismem 1968 osób. Najlepszy jest Qwen2.5-VL 7B (CER 0,102, WER 0,242),
następnie InternVL3-8B (0,219) i~DeepSeek-VL2 (0,355). Silniki klasyczne wypadają gorzej:
Tesseract 0,419, EasyOCR 0,632 i~PaddleOCR 0,649, a~GOT-OCR2 prawie nie rozpoznaje pisma (0,795).
PaddleOCR zwraca pustą odpowiedź dla 17 skanów, Tesseract dla 3, EasyOCR dla 2.

\begin{table}
  \caption{Średni CER na zdaniach pisma odręcznego (syntetyczne, obrazy czyste) i~na skanach Gold Setu,
  miejsce w~rankingu oraz WER i~liczba pustych odpowiedzi na skanach.}
  \label{tab:gold}
  \centering\small
  \begin{tabular}{lrrrrrr}
\toprule
Model & \multicolumn{2}{c}{CER} & \multicolumn{2}{c}{Miejsce} & WER & Puste \\
 & syntetyczny & PHSF & syntetyczny & PHSF & PHSF & PHSF \\
\midrule
Qwen2.5-VL 7B & .012 & .103 & 1 & 1 & .242 & 0 \\
InternVL3-8B & .087 & .220 & 2 & 2 & .499 & 1 \\
DeepSeek-VL2 & .101 & .355 & 3 & 3 & .712 & 0 \\
Tesseract & .164 & .419 & 5 & 4 & .862 & 3 \\
EasyOCR & .207 & .632 & 6 & 5 & .989 & 2 \\
PaddleOCR & .131 & .650 & 4 & 6 & .933 & 17 \\
GOT-OCR2 & .328 & .795 & 7 & 7 & .982 & 0 \\
\bottomrule
\end{tabular}

\end{table}

\textbf{Zgodność rankingów (pytanie 3).} Korelacja Spearmana między rankingiem na syntetycznym piśmie
odręcznym a~rankingiem na skanach wynosi $\rho=0{,}89$ ($p=0{,}007$, $n=7$), czyli przekracza próg
istotności dla siedmiu modeli ($\rho\approx0{,}79$ przy poziomie 0,05, test dwustronny). Trzy najlepsze
i~najgorszy model zajmują te same miejsca w~obu rankingach; różnią się tylko miejsca Tesseracta, EasyOCR
i~PaddleOCR. Dla porównania autorzy badania bazowego podają $\rho=0{,}92$.

%------------------------------------------------------------------------------
\section{Dyskusja}

\textbf{Ranking i~rodzaj modelu.} Na obu zbiorach najlepsze są ogólne modele wizyjno-językowe
(Qwen, InternVL3, DeepSeek-VL2), a~najlepszy z~nich, Qwen2.5-VL, jest wyraźnie lepszy od reszty. Klasyczne
silniki z~modelem polskim są najlepsze na czystym druku, lecz nie radzą sobie z~pismem odręcznym
i~zniekształceniami. GOT-OCR2, choć wyspecjalizowany w~OCR, jest ostatni na obu zbiorach.

\textbf{Słowniki znaków.} PaddleOCR dla języka polskiego korzysta ze wspólnego modelu łacińskiego,
którego słownik (\texttt{latin\_dict.txt}) zawiera litery ą, ć, ę, ł, ó, ale nie zawiera ń, ś, ź, ż.
Model nie może więc ich zwrócić i~zastępuje je literami podobnymi (np. ,,ludnošci'' zamiast
,,ludności''). 

\textbf{Odtwarzalność badania bazowego.} Ponowne uruchomienie ośmiu otwartych modeli z~badania bazowego
na próbie 1650 obrazów dało średni CER różniący się od opublikowanych predykcji najwyżej o~0,001
(Pixtral: 0,007). Kod autorów liczy jednak CER, dzieląc przez długość predykcji, a~nie wzorca, a~liczby
z~tekstu artykułu odtwarza inny wariant z~ich repozytorium (bez spacji i~interpunkcji), czego artykuł
nie opisuje. Zmienia to ranking: GOT-OCR2, który gubi spacje, ma w~tym wariancie CER 0,311, a~wzorem~\eqref{eq:cer} 0,391.

\TODO{Porównanie z~wynikami tureckimi (tylko jakościowo, bez zestawiania liczb).}

%------------------------------------------------------------------------------
\section{Ograniczenia i dalsze prace}

Zbiór syntetyczny jest jednoliniowy i~na białym tle (tekstury papieru z~badania bazowego nie zostały
udostępnione), a~pominięto tekst w~scenie. Skany obejmują jeden krótki tekst. Do porównania weszło siedem modeli, a~nie dziewięć: Phi-3.5-vision i~Pixtral-12B
pozostały niedokończone. Przy $n=7$ korelacja Spearmana jest istotna tylko dla dużych wartości
($\rho\gtrsim0{,}79$), a~jej wartość zależy od rankingu trzech środkowych modeli. Wyniki pochodzą
z~jednego przebiegu (ziarno 42).

%------------------------------------------------------------------------------
\section{Wnioski}

W~naszym badaniu Qwen2.5-VL 7B jest najdokładniejszy zarówno
na danych syntetycznych (CER 0,059), jak i~na skanach odręcznego pisma (0,102). Polskie znaki pogarszają
wyniki większości modeli, najbardziej GOT-OCR2, DeepSeek-VL2, PaddleOCR i~InternVL3, a~klasyczne silniki
są wrażliwe na szum, rozmycie i~niską rozdzielczość. Ranking na syntetycznym piśmie odręcznym odpowiada
rankingowi na rzeczywistych skanach ($\rho=0{,}89$), więc tanie dane syntetyczne dobrze przewidują kolejność
modeli na prawdziwym piśmie, choć nie jego poziom błędu.

\subsubsection*{Dostępność danych i kodu.}
Generator, skrypty uruchomieniowe, wyniki dla każdego obrazu i~zbiór obrazów:
\TODO{adres repozytorium}. Kod na licencji MIT, dane na licencji CC BY 4.0.

\bibliographystyle{splncs04}
\bibliography{references}

\end{document}
